\documentclass[11pt]{article}
\usepackage[margin=1in]{geometry}
\usepackage{booktabs,graphicx,hyperref,amsmath}
\title{Fine-tuning a 30B Language Model to Write Low-Poly 3D Meshes as Text:\\ What Scale and Canonical Ordering Buy on a Small Budget}
\author{Irfan Khan}
\date{DRAFT \today}
\begin{document}
\maketitle

\begin{abstract}
We fine-tune Muse Glimmer 30B (Apache 2.0) with LoRA to generate small coloured triangle meshes as plain text
(\texttt{v x y z} on a $64^3$ grid, \texttt{c \#rrggbb} colour groups, \texttt{f a b c} faces), in the style of LLaMA-Mesh.
Across four training runs on a single rented A100 (total GPU spend in the tens of dollars), the share of held-out captions that yield a
syntactically valid mesh rises from 17\% (73 training objects) to 80\% (1,298 objects with canonical vertex ordering);
the share of valid meshes that a human judges recognisably correct rises from about 1 in 6 to about 1 in 3.
\textbf{[FINAL RUN: N objects, valid X\%, recognisable Y\% --- to fill in.]}
We report the failure modes (runaway coordinate ``counting'' loops, right colour and size but wrong detail), a decoding
guard that removes the loops, and a licensing-aware data pipeline (CC0 and CC-BY only, with an attribution manifest).
All numbers are on small held-out sets (30 captions) and the data-scale and ordering changes are confounded.
\end{abstract}

\section{Introduction}
Language models can emit meshes as text, which makes 3D generation a fine-tuning problem rather than an architecture
problem \cite{llamamesh}. We ask how far a modest LoRA run gets with a large general model, little compute and openly licensed data.

\section{Method}
\paragraph{Mesh format.} Meshes are normalised to a unit cube, quantised to 0--63, deduplicated and written as text.
Optionally (``sorted'') vertices are ordered by $(z,y,x)$ and each colour group's faces by lowest vertex index, with each face rotated
so the lowest index comes first, following \cite{llamamesh}. Sorting is re-applied after every rotation or mirror augmentation.
\paragraph{Training.} Muse Glimmer 30B, BF16 LoRA (rank 16, attention and MLP projections of the language model), lr $10^{-4}$, 2 epochs,
one A100 80GB, loss on the mesh completion only. Prompts are a short system message plus a caption.
\paragraph{Decoding.} Four samples per caption at temperature 0.6. A loop guard stops a sample when its last six vertex lines step by the same nonzero amount
(a runaway ``counting'' loop). The returned sample is the first valid one, or (final run) the medoid of valid samples in size, triangle count and colour count.
\paragraph{Validity.} A mesh is valid if it parses, indices are in range, there are enough vertices and faces, it is not flat, spans enough of the grid and has no unused vertices.

\section{Data}
\begin{itemize}
\item \textbf{Hand-captioned CC0 set (166 $\to$ 126 after removing unverifiable sources):} Quaternius and Kenney-style kits converted from GLB, 69 of them Blender-decimated to 90--150 triangles, captions written from renders.
\item \textbf{Objaverse CC-BY subset:} from the 798k Objaverse annotations, licence \texttt{by}, 40--300 faces, no animation, useful categories. Captions from Cap3D \cite{cap3d}.
Filtered to $\le$250 lines and to drop captions naming a colour when the mesh is grey (Objaverse colours live in textures we do not read).
All Objaverse CC0 objects were high-detail scans and unusable. 887 distinct authors in the first 1,278 objects; attribution manifest released.
\item Excluded on licence grounds: Kenney (site declares an AI-training opt-out), Quaternius Asset License v1.0 kits, NC and SA Objaverse objects.
\end{itemize}

\section{Results}
\begin{table}[h]\centering
\begin{tabular}{lrrrr}
\toprule
Run & Train objects & Held-out loss (before$\to$after) & Valid / captions & Recognisable (of valid) \\
\midrule
A & 73 & 0.607$\to$0.431 & 3/18 (17\%) & -- \\
A + guard, 4 samples & 73 & -- & 6/18 (33\%) & 1 of 6 \\
B & 130 & 0.696$\to$0.508 & 11/30 (37\%) & about 1 in 6 \\
C (sorted) & 1{,}298 & 0.514$\to$0.217 & 24/30 (80\%) & about 1 in 3 \\
D (final, 3,497 sorted, 45 prompts) & 0.216 held-out loss & 34/45 (76\%); 21/30 on the run-C prompts (C: 24/30) & about 15--20\% recognisable (visual) & no gain over C \\
\bottomrule
\end{tabular}
\caption{Held-out sets differ between runs, so losses are not directly comparable. ``Recognisable'' is the author's visual judgement of rendered output against the target.}
\end{table}
\paragraph{Failure modes.} (i) Runaway arithmetic loops (\texttt{v 31 63 32, v 31 63 31, \dots}) until the token cap; 13 of 18 captions in run A, 9 of 30 in run B, 3 of 30 in run C.
(ii) Valid meshes with the right colour and bounding box but wrong detail (apple $\to$ flat disc, houseplant $\to$ cone, wireframe prism $\to$ solid box).
(iii) Qwen 0.5B (earlier rung): held-out loss rose (0.73$\to$1.95) with 17 training objects, i.e.\ memorisation.

\section{Limitations}
Thirty held-out captions per run; recognisability is one rater's judgement; run C changed object count, mesh length and ordering together, so we cannot attribute the gain; meshes are small (up to about 250 lines); Objaverse colours from textures are lost, so many objects are grey; per-sample validity is lower than per-caption validity because we keep the best of four.

\section{Release}
LoRA adapter, evaluation outputs, scripts, and \texttt{objaverse\_attribution.csv} (uid, author, licence, URL for every CC-BY object). The base model is Apache 2.0; CC-BY authors are credited in the manifest. \textbf{[Licence of the adapter: to decide.]}

\begin{thebibliography}{9}
\bibitem{llamamesh} Z.~Wang et al. LLaMA-Mesh: Unifying 3D Mesh Generation with Language Models. arXiv:2411.09595, 2024.
\bibitem{cap3d} T.~Luo et al. Scalable 3D Captioning with Pretrained Models (Cap3D). \textbf{[verify authors and venue before submission]}
\bibitem{objaverse} M.~Deitke et al. Objaverse: A Universe of Annotated 3D Objects. \textbf{[verify]}
\end{thebibliography}
\end{document}
